\section*{Materials and Methods}

\subsection*{The \AP{} Hidden Markov Model}

\paragraph{Overview.}
\AP{} infers archaic introgression at haplotype resolution by extending the
Li \& Stephens (2003) haplotype copying model~\citep{li2003modeling} to include
archaic reference genomes as competing ancestry sources.
The core insight is coalescent-theoretic: a modern human haplotype segment that traces
ancestry to an archaic hominin via introgression coalesces with the archaic reference
at time $t_\mathrm{admixture} \approx 1{,}724$ generations ($\sim$50 kya), whereas a
non-introgressed segment coalesces at the modern human--archaic split time
$t_\mathrm{split} \approx 19{,}000$ generations ($\sim$550 kya).
Under a Poisson recombination model, the expected tract length scales as $1/t$ Morgans,
yielding an 11-fold difference in expected haplotype match lengths
($\approx$58 kb for introgressed versus $\approx$5 kb for non-introgressed haplotypes
at a typical recombination rate of 1~cM/Mb).
\AP{} exploits this \emph{match-length signal} via exponential transition
probabilities rather than the allele density signal used by existing tools.
Methods such as HMMix~\citep{skov2018detecting}, DAIseg~\citep{planche2024daiseg},
and IBDmix~\citep{chen2020identifying} model archaic introgression via Poisson
counts of archaic-specific alleles within sliding windows, an architecture that
discards all haplotype linkage disequilibrium structure.
The tract-length transition formulation is analogous to modern local ancestry
inference (HAPMIX~\citep{price2009sensitive};
MOSAIC~\citep{salter2019mosaic}; ChromoPainter~\citep{lawson2012inference}),
here adapted to archaic reference genomes.

\paragraph{Hidden states.}
\AP{} assigns each phased SNP position on a query haplotype to one of three
mutually exclusive hidden states: $S_\mathrm{AMH}$ (modern human ancestry),
$S_\mathrm{NEA}$ (Neanderthal-introgressed), and $S_\mathrm{DEN}$
(Denisovan-introgressed).
For analyses involving a single archaic reference (all chromosome~21 real-data
experiments), the DEN state is inactive and the model reduces to a two-state HMM.
Initial state probabilities are set to $\pi_\mathrm{NEA} = 0.02$ and
$\pi_\mathrm{AMH} = 0.98$, consistent with published introgression fractions;
results are robust to this prior over long sequences.

\paragraph{Emission probabilities.}
At each biallelic SNP site $i$, the query haplotype carries allele $a_i \in \{0, 1\}$.
Let $f_i^\mathrm{YRI}$ denote the derived allele frequency estimated from Yoruban
(YRI) haplotypes in the 1000~Genomes Project Phase~3 panel~\citep{genomes20151000}.
Under $S_\mathrm{AMH}$:
\begin{equation}
  P(a_i \mid S_\mathrm{AMH}) = f_i^\mathrm{YRI} \cdot [a_i=1] + (1 - f_i^\mathrm{YRI}) \cdot [a_i=0].
  \label{eq:emiss_amh}
\end{equation}

Under $S_\mathrm{NEA}$, two emission models are used depending on whether the
archaic reference genome matches the true introgressing population.

\emph{Standard bidirectional emission (simulation).}
When the reference is drawn from the same simulated population as the true donor,
mismatches are informative and the standard Li \& Stephens emission is used:
\begin{equation}
  P(a_i \mid S_\mathrm{NEA}, \text{sim.}) =
  (1-\theta)[a_i = g_i^\mathrm{arch}] + \theta[a_i \neq g_i^\mathrm{arch}],
  \label{eq:emiss_nea_std}
\end{equation}
where $\theta = 0.01$ and $g_i^\mathrm{arch}$ is the archaic reference allele.
On the 50-replicate benchmark this formulation achieves F1 $= 0.480 \pm 0.095$.

\emph{Positive-only emission (real data).}
Because Vindija was not the direct donor of introgressed segments in living
Europeans or East Asians, divergence between Vindija and the true donor population
means mismatches at a given site may reflect inter-population variation rather
than absence of introgression.
We therefore treat mismatches as uninformative for real-data analyses:
\begin{equation}
  P(a_i \mid S_\mathrm{NEA}, \text{real}) =
  \begin{cases}
    1 - \theta & \text{if } a_i \text{ matches } g_i^\mathrm{arch}, \\
    \tfrac{1}{2}        & \text{otherwise.}
  \end{cases}
  \label{eq:emiss_nea_pos}
\end{equation}
Setting $P(\text{mismatch} \mid S_\mathrm{NEA}) = 1/2$ assigns zero log-odds
to disagreements, equivalent to treating mismatches as uninformative when the
reference diverges from the true donor.
Critically, the standard and positive-only emission models are \emph{not}
interchangeable: on simulation, the positive-only model achieves F1 $= 0.000$
because mismatches carry essential discriminative information when the reference
exactly matches the source.
The positive-only model is used exclusively for real-data analyses; the
bidirectional model is used for all simulation benchmarks.

\paragraph{Transition probabilities.}
Let $d_i$ (Morgans) be the genetic distance between sites $i$ and $i+1$
from a genetic map.
Under the Haldane recombination model~\citep{price2009sensitive}:
\begin{equation}
  P(S_{i+1} = s \mid S_i = s) = e^{-d_i \cdot t_s},
  \label{eq:trans}
\end{equation}
where $t_s$ is the time parameter in generations:
$t_\mathrm{NEA} = 1{,}724$, $t_\mathrm{DEN} = 1{,}000$,
$t_\mathrm{AMH} = 19{,}000$.
The expected tract length is $1/t_s$ Morgans
($\approx$58 kb for NEA; $\approx$100 kb for DEN; $\approx$5 kb for AMH
at 1~cM/Mb).
Off-diagonal probabilities are distributed uniformly.
Genetic distances are computed by linear interpolation of the GRCh38
sex-averaged recombination map~\citep{halldorsson2019characterizing}
(HapMap-II format, UCSC Genome Browser).

\paragraph{Viterbi decoding and post-processing.}
State sequences are inferred by the Viterbi algorithm in log-space.
Consecutive archaic-state positions are merged into segments.
For real-data analyses, gaps $\leq 10$~kb between adjacent segments are closed
(reducing fragmentation from sparse Vindija sites) and merged segments shorter
than 10~kb are discarded.
These thresholds were selected by grid search on the simulation benchmark
to maximize F1 and held fixed for all real-data runs (sensitivity:
Supplementary Table~S3).

\subsection*{Simulation Experiments}

Coalescent simulations used \texttt{msprime}~v1.3.3~\citep{baumdicker2022efficient}
(Python~3.12).
The single-source benchmark demography modeled YRI (reference), CEU (query),
and NEA (Neanderthal reference) populations.
A population split at $t_\mathrm{OOA} = 2{,}000$ generations separated YRI and CEU,
followed by a 2\% admixture pulse from NEA into CEU at $t_\mathrm{admixture} = 1{,}724$
generations via \texttt{msprime.MassMigration}.
The NEA--modern split was at $t_\mathrm{split} = 19{,}000$ generations~\citep{prufer2014complete};
all $N_e = 10{,}000$ diploid.

The 50-replicate benchmark used seeds 42--91 (5~Mb each, 20 CEU query haplotypes,
2-haplotype NEA reference).
Ground-truth intervals were extracted from all ARG migration records with times
within $\pm 80\%$ of $t_\mathrm{admixture}$.

The two-source proof-of-concept used a two-archaic demography with
MassMigration events at $t_\mathrm{NEA} = 1{,}724$ gen (2\%) and $t_\mathrm{DEN} = 1{,}000$
gen (1\%), and population splits at $t_\mathrm{split\_NEA} = 19{,}000$ gen and
$t_\mathrm{split\_DEN} = 25{,}000$ gen.
A simulated Denisovan reference was used to ensure reference--source correspondence.
Twenty replicates were run (seeds 100--119, 5~Mb, 20 query haplotypes).

\subsection*{Evaluation Metrics}

\emph{F1}: true positive = predicted segment overlaps $\geq 1$~bp of a ground-truth segment
(lenient, following Sankararaman et al.~\citeyear{sankararaman2014genomic};
stricter 50\% and 80\% reciprocal-overlap results in Supplementary Table~S4).
\emph{AUPRC}: precision-recall curve integrated over posterior probability thresholds.
\emph{Attribution accuracy}: fraction of predicted NEA (or DEN) segments overlapping
a same-type ground-truth segment; \emph{cross-contamination}: fraction overlapping
the opposite-type ground truth.
\emph{Bootstrap CIs}: percentile bootstrap, 2{,}000 resamples (SciPy~1.14.1~\citep{virtanen2020scipy}).
\emph{Wilcoxon test}: one-sided signed-rank test for pairwise method comparisons.

\subsection*{Real-Data Analysis}

\paragraph{Data.}
Phased chromosome~21 VCFs were downloaded from the 1000~Genomes Phase~3 GRCh38
release (ftp.1000genomes.ebi.ac.uk)~\citep{genomes20151000}.
We analyzed 20 CEU and 20 CHB individuals (40 phased haplotypes each)
across chr21:10,000,000--48,000,000 (38~Mb).
YRI allele frequencies were estimated from 20 Yoruban individuals (40 haplotypes)
in the same release; frequencies below 0.025 were floored at 0.025 to prevent
emission instability from singletons.

\paragraph{Vindija reference.}
The Vindija~33.19 diploid VCF~\citep{prufer2017high} was downloaded from the
MPI EVAN repository, converted from GRCh37 to GRCh38 via CrossMap~v0.7.0~\citep{zhao2014crossmap},
and merged with the 1000GP VCF using allele-aware orientation correction
(swapped REF/ALT codes were inverted before computing match scores).
After intersection and quality filtering, 1{,}036 (CEU) and 906 (CHB) informative
SNPs remained ($\approx$27 SNPs/Mb), substantially fewer than in simulation
($\approx$320 SNPs/Mb) due to the limited callable regions of the Vindija genome.
The median inter-SNP distance ($\approx$37~kb) retains sensitivity to tracts
$\geq 50$~kb but reduces boundary resolution relative to simulation.

\paragraph{Ancestry fraction.}
Per-haplotype Neanderthal ancestry fraction is the proportion of the 38-Mb window
covered by NEA-state Viterbi calls.
Population estimates are mean $\pm$ SD across the 40 haplotypes.

\subsection*{Two-Source Attribution}
\label{sec:attribution}

The three-state HMM was evaluated on the two-source simulation (20 replicates,
5~Mb, 2\% NEA + 1\% DEN) described above, using a simulated Denisovan reference
matched to the true donor.
Real-data Denisovan attribution using the Altai genome is deferred to future work
owing to the absence of a GRCh38-compatible high-coverage Denisovan VCF for chr21.

\subsection*{Functional Annotation}

Introgressed segment coordinates were intersected with UCSC knownGene
annotations (hg38, downloaded February~2026)~\citep{kent2002human};
gene symbols were mapped via the UCSC kgXref table.
A transcript was considered overlapping at $\geq 1$~bp.
Named genes were cross-referenced against a curated list of chr21 adaptive
introgression loci (Supplementary Table~S2) compiled from Vernot \& Akey~\citeyear{vernot2014resurrecting},
Sankararaman et al.~\citeyear{sankararaman2014genomic},
Vernot et al.~\citeyear{vernot2016excavating},
Browning et al.~\citeyear{browning2018ancestry},
and Dannemann \& Kelso~\citeyear{dannemann2017contribution}.

\subsection*{Scalability Analysis}

Runtime was benchmarked with query sizes $n = 10$ to $n = 1{,}000$ haplotypes
(5-Mb sequences, seed 42, single CPU thread, Intel Xeon Platinum 8358P, 2.60~GHz).
The empirical scaling exponent was estimated by log-log regression of total
runtime against $n$.
Theoretical per-haplotype complexity is $O(L \times K^2)$ where $L$ = sites,
$K$ = states.

\subsection*{Software and Reproducibility}

\AP{} is implemented in Python~3.12; code, simulation scripts, preprocessing
pipelines, and pinned dependencies (\texttt{environment.yml}) are provided
in Supplementary Scripts~S1--S5 and at \url{https://github.com/[repository]}.
